% TOP_PROBLEM: {"id": 30006753, "title": "Global Regularity of Solutions to the Inviscid Surface Quasi-Geostrophic Equation", "category_id": 9, "category": {"id": 9, "name": "pde", "display_name": "Partial Differential Equations", "description": "PDEs and their applications in physics and geometry.", "slug": "pde", "order_index": 9, "created_at": "2026-07-31T15:26:25.671Z"}, "status": "open"}
% FIRST_POSTED: 2026-09-25
% ENTRY_KIND: statement_only
% !TeX program = lualatex
% Definition and sources only; this file is not an LLM solution attempt.
\documentclass[11pt]{article}
\usepackage[margin=25mm]{geometry}
\usepackage{fontspec}
\setmainfont{Latin Modern Roman}
\usepackage{amsmath,amssymb,mathtools}
\usepackage{unicode-math}
\setmathfont{Latin Modern Math}
\usepackage{xurl,hyperref}
\hypersetup{colorlinks=true,urlcolor=blue}
\setcounter{secnumdepth}{0}
\setlength{\parindent}{0pt}
\setlength{\parskip}{5pt}
\setlength{\emergencystretch}{3em}
\begin{document}
\section*{\texorpdfstring{Global Regularity of Solutions to the Inviscid Surface Quasi-Geostrophic Equation}{Global Regularity of Solutions to the Inviscid Surface Quasi-Geostrophic Equation}}\label{30006753}
\subsection{Definitions and mathematical statement}
For the scalar $\theta(t,x)$ on ${\mathbb{R}}^2$, let $R_j$ be the Riesz transform with Fourier multiplier $-i\xi_j/|\xi|$ and set $R^\perp\theta=(-R_2\theta,R_1\theta)$. Inviscid SQG is
\[
 \partial_t\theta+u\cdot\nabla\theta=0,\qquad u=R^\perp\theta,
 \qquad\theta(0,x)=\theta_0(x).
\]
The problem is global smooth existence for every permitted smooth decaying finite-energy datum, or a finite-time breakdown for one such datum. ``Inviscid'' means no dissipative fractional-Laplacian term is present. The precise initial norm and smoothness class should follow the cited formulation rather than a weaker low-regularity ill-posedness result.
\par\smallskip\textit{Formulation and concept references:} \href{https://www.sciencedirect.com/science/article/pii/S0001870825004517}{[S1]}.

\subsection{Short English statement}
Can a smooth inviscid SQG flow form a singularity in finite time, or must every allowed smooth initial state remain smooth forever?
\subsection{Sources}
\begin{itemize}
\item[S1]Cordoba, Martinez-Zoroa, and Ozanski, Instantaneous Continuous Loss of Regularity for the SQG Equation, Introduction.
\url{https://www.sciencedirect.com/science/article/pii/S0001870825004517}.
\textit{Formulation source.}
\end{itemize}
\end{document}
